% Torque-free symmetric top seen in the body frame (oblate case I3 > I1, omega > 0).
% Omega = (A cos(wt), A sin(wt), Omega3) sweeps a cone about x3; M = (I1 Omega1, I1 Omega2, I3 Omega3)
% lies in the plane of x3 and Omega and sweeps a narrower cone (drawn with I3/I1 = 1.5).
% Compile: xelatex symmetric_top_free_rotation_tikz.tex ; pdftoppm -png -r 300 -singlefile symmetric_top_free_rotation_tikz.pdf symmetric_top_free_rotation
\documentclass[tikz,border=4pt]{standalone}
\usepackage{amsmath}
\usepackage[fontset=fandol]{ctex}
\usepackage{tikz-3dplot}
\usetikzlibrary{arrows.meta}

\begin{document}
\tdplotsetmaincoords{62}{110}
\begin{tikzpicture}[tdplot_main_coords,>=Stealth,line cap=round,
                    every node/.style={font=\normalsize}]
  \def\A{1.35}\def\H{2.6}\def\k{1.5}      % |Omega_perp|, Omega_3, I3/I1
  \def\ta{100}                              % azimuth of Omega at the instant drawn
  % M/|M| scaled to length 3.0
  \pgfmathsetmacro{\mn}{sqrt(\A*\A+\k*\k*\H*\H)}
  \pgfmathsetmacro{\mr}{3.0*\A/\mn}\pgfmathsetmacro{\mh}{3.0*\k*\H/\mn}

  % body axes
  \draw[->,gray!70!black] (0,0,0) -- (2.0,0,0) node[anchor=north east] {$x_1$};
  \draw[->,gray!70!black] (0,0,0) -- (0,2.0,0) node[anchor=west] {$x_2$};
  \draw[->,blue!75!black,thick] (0,0,0) -- (0,0,3.6) node[anchor=south] {$x_3$（对称轴）};

  % cone of Omega (body cone) and cone of M
  \fill[orange!12,opacity=.7] plot[domain=0:360,samples=80,variable=\t] ({\A*cos(\t)},{\A*sin(\t)},\H) -- cycle;
  \draw[orange!85!black] plot[domain=0:360,samples=80,variable=\t] ({\A*cos(\t)},{\A*sin(\t)},\H);
  \draw[orange!50,thin] (0,0,0) -- ({\A*cos(110)},{\A*sin(110)},\H);
  \draw[orange!50,thin] (0,0,0) -- ({\A*cos(-70)},{\A*sin(-70)},\H);
  \draw[green!45!black,densely dashed] plot[domain=0:360,samples=80,variable=\t] ({\mr*cos(\t)},{\mr*sin(\t)},\mh);

  % sense of rotation (counter-clockwise seen from +x3 when omega > 0)
  \draw[->,orange!85!black,very thick] plot[domain=-40:25,samples=20,variable=\t] ({\A*cos(\t)},{\A*sin(\t)},\H);
  \node[orange!85!black,anchor=north] at ({\A*cos(-10)},{\A*sin(-10)},\H) {$\omega$};

  % vectors
  \draw[densely dashed] (0,0,0) -- (0,0,\H);
  \draw[->,green!45!black,very thick] (0,0,0) -- ({\mr*cos(\ta)},{\mr*sin(\ta)},\mh) node[pos=.72,anchor=east] {$\boldsymbol M$};
  \draw[->,red,very thick] (0,0,0) -- ({\A*cos(\ta)},{\A*sin(\ta)},\H) node[anchor=west] {$\boldsymbol\Omega$};
  \draw[->,purple,thick] (0,0,\H) -- ({\A*cos(\ta)},{\A*sin(\ta)},\H) node[pos=.55,anchor=south] {$\boldsymbol\Omega_\perp$};
  \node[anchor=east] at (0,0,{0.45*\H}) {$\Omega_3$};
  \fill (0,0,0) circle (1.4pt);

  \begin{scope}[tdplot_screen_coords]
    \node[anchor=west,align=left] at (2.9,1.6)
      {$\Omega_1=A\cos\omega t,\quad \Omega_2=A\sin\omega t$\\[3pt]
       $\Omega_3=\text{常数},\quad \omega=\dfrac{I_3-I_1}{I_1}\Omega_3$\\[8pt]
       {\small 刚体系中：$\boldsymbol\Omega$ 与 $\boldsymbol M$ 都绕对称轴}\\
       {\small 以角速度 $\omega$ 转动，三者始终共面}\\[2pt]
       {\small （图中取扁陀螺 $I_3=1.5I_1$，$\omega>0$）}};
  \end{scope}
\end{tikzpicture}
\end{document}
